% Ad Atticum 3.9 (ad-atticum-03-09)
% Source: https://raw.githubusercontent.com/PerseusDL/canonical-latinLit/master/data/phi0474/phi057/phi0474.phi057.perseus-lat1.xml
% Fetched by scripts/fetch_latin.py

% Dateline (Perseus): Scr. Thessalonicae Id. Iun. a. 696 (58) .

\ciceroLetterOpener{CICERO ATTICO salutem}

\ciceroSection{1} Quintus frater cum ex Asia discessisset ante Kal. Maias et Athenas venisset Idibus, valde fuit ei properandum, ne quid absens acciperet calamitatis, si quis forte fuisset qui contentus nostris malis non esset. itaque eum malui properare Romam quam ad me venire et simul (dicam enim quod verum est, ex quo magnitudinem miseriarum mearum perspicere possis) animum inducere non potui ut aut illum amantissimum mei mollissimo animo tanto in maerore aspicerem aut meas miserias luctu adflictus et perditam fortunam illi offerrem aut ab illo aspici paterer. atque etiam illud timebam, quod profecto accidisset, ne a me digredi non posset. versabatur mihi tempus illud ante oculos quom ille aut lictores dimitteret aut vi avelleretur ex complexu meo. huius acerbitatis eventum altera acerbitate non videndi fratris vitavi. in hunc me casum vos vivendi auctores impulistis.

\ciceroSection{2} itaque mei peccati luo poenas. quam quam me tuae litterae sustentant ex quibus quantum tu ipse speres facile perspicio; quae quidem tamen aliquid habebant solaci ante quam eo venisti a Pompeio. nunc Hortensium adlice et eius modi viros. obsecro , mi Pomponi, nondum perspicis quorum opera, quorum insidiis, quorum scelere perierimus? sed tecum haec omnia coram agemus; tantum dico quod scire te puto, nos non inimici sed invidi perdiderunt. nunc si ita sunt quae speras, sustinebimus nos et spe qua iubes nitemur; sin, ut mihi videntur, infirma sunt, quod optimo tempore facere non licuit minus idoneo fiet.

\ciceroSection{3} Terentia tibi saepe agit gratias. mihi etiam unum de malis in metu est, fratris miseri negotium; quod si sciam quoius modi sit, sciam quid agendum mihi sit. me etiam nunc istorum beneficiorum et litterarum exspectatio, ut tibi placet, Thessalonicae tenet. si quid erit novi adlatum, sciam de reliquo quid agendum sit. tu si, ut scribis, Kal. Iuniis Roma profectus es, prope diem nos videbis. litteras quas ad Pompeium scripsi tibi misi. data id . Iun. Thessalonicae.
